\documentclass[10pt,a4paper]{amsart}
\usepackage[margin=19mm]{geometry}
\usepackage{amsmath,amssymb,mathtools}
\usepackage[scr=rsfs,cal=euler]{mathalfa}
\usepackage{xfrac}
\usepackage[unicode,hidelinks,bookmarks=false]{hyperref}
\newcommand{\ie}{i.e.\ }
\newcommand{\cf}{cf.\ }
\newcommand{\resp}{respectively\ }
\newcommand{\Subsection}{\subsection}
\providecommand{\nobreakdash}{\nobreak\hbox{-}}
\setlength{\emergencystretch}{2em}
\multlinegap0pt
\usepackage[shortlabels]{enumitem}
\usepackage{tikz}
\usepackage{amssymb}
\usepackage{siunitx}
\usepackage{multirow}
\newtheorem{prop}{Proposition}
\title{Combinatorial invariants for certain classes of non-abelian groups}
\author{Naveen K. Godara, Renu Joshi,, Eshita Mazumdar}
\date{Source: arXiv:2408.13558v2 (CC BY 4.0)}
\begin{document}
\maketitle

\noindent\emph{Source and licence.} Combinatorial invariants for certain classes of non-abelian groups. Version arXiv:2408.13558v2 (CC BY 4.0), distributed by arXiv under the Creative Commons Attribution 4.0 International licence, http://creativecommons.org/licenses/by/4.0/, as recorded in the arXiv licence metadata for this exact version and shown on the abstract page https://arxiv.org/abs/2408.13558v2. Full licence notice: \texttt{licenses/arxiv-2408.13558-CC-BY-4.0.txt}.\par\smallskip

\noindent\textit{Editorial scope.} Counted unit: the article's prop Ind (source0.tex lines 299--304). The article's complete original proof of it is delivered with it (source0.tex lines 305--328, one \texttt{proof} environment of 24 lines). No mathematics is written, repaired or re-derived: the statement and the proof are the article's own lines, in the article's own notation, and no retained line is substituted or re-flowed.  Retained uncounted as the source-local context the proof needs: lines 294--296.  Nothing was omitted from any retained span, so the package carries no omission record.  No retained line contains a citation, so the wrapper carries no bibliography and no bibliography entry.  No retained line contains a figure environment, an image-inclusion command or drawing code, so no image is delivered and the package declares no asset dependency.  Editorial formatting only, in addition: an AMS \texttt{amsart} preamble that restates the article's own declarations and macros used by the retained lines, namely the environments \texttt{prop} on separate counters, together with the standard packages those lines need.  The retained environments print in this document as Proposition 1 in order of appearance, whereas the article prints the same items with the numbering of its own full document.  The complete native source of the article as distributed in the arXiv e-print for this exact version is bundled unchanged as \texttt{sources/164218/source0.tex}.
\begin{equation}
[G,G] \neq 1 \quad \text{and}\quad [[G,G],G]=1. \label{2.1}
\end{equation}

\begin{prop}\label{Ind} 
Let $G$ be a $2$-group satisfying relation \emph{(\ref{2.1})}. Then, for all $s \geq 1$, we have 
\begin{equation*}\label{equationinduction}
M_j=G^{2^{s}} \;\text{ for all }\; j \in [2^{s-1}+1, 2^s].
\end{equation*}
\end{prop}

\begin{proof} Assume that
$P(s) :  M_{2^{s-1}+1}= \cdots =M_{2^{s}}=G^{2^{s}} \; \text{for} \; s \ge 1.$ We will use induction on $s \ge 1$. From the definition of the $M$-series of $G$, we have $M_2=[G,G]G^2$. Since $[G,G]\subseteq G^2$, it follows that $P(1)$ holds. Assume that $P(s)$ is true for some $s\ge 1$, i.e.
\begin{equation}\label{eq:indhyp}
M_{2^{s-1}+1}=\cdots =M_{2^{s}}=G^{2^{s}}.
\end{equation}
We show that $P(s+1): M_{2^{s}+1}=\cdots=M_{2^{s+1}}=G^{2^{s+1}}$ holds. First, using the recursive definition 
\[
M_{2^s+1}
=[M_{2^s},G]\, M_{\lceil (2^s+1)/2\rceil}^{2},
\]
 and since $\lceil (2^{s}+1)/2\rceil=2^{s-1}+1$, the induction hypothesis
\eqref{eq:indhyp} gives $M_{2^{s}}=G^{2^{s}}$ and $M_{2^{s-1}+1}=G^{2^{s}}$. Hence
$M_{2^{s}+1}
=[G^{2^{s}},G]\,(G^{2^{s}})^{2}
=G^{2^{s+1}},$ because 
$[G^{2^{s}},G]\subseteq G^{2^{s+1}}$. Now, assume that $$M_t=G^{2^{s+1}} \text{ for all } t \in [2^{s}+1, k]\quad \text{for some $k$ with $2^{s}+1\le k<2^{s+1}$}.$$ Then \[
M_{k+1}=[M_k,G]\; M_{\lceil (k+1)/2\rceil}^{2} = [G^{2^{s+1}},G]\; M_{\lceil (k+1)/2\rceil}^{2}.
\]
Since $k+1\le 2^{s+1}$, it follows that $\lceil (k+1)/2\rceil\in[2^{s-1}+1,2^{s}]$, and the induction hypothesis \eqref{eq:indhyp} implies $M_{\lceil (k+1)/2\rceil}=G^{2^{s}}$. Therefore
\[
M_{k+1} = [G^{2^{s+1}},G]\,(G^{2^{s}})^2 = G^{2^{s+1}}, \text{ since $[G^{2^{s+1}},G]=[G,G]^{2^{s+1}}\subseteq G^{2^{s+2}}\subseteq G^{2^{s+1}}.$}
\] By induction on $t$, we conclude that 
$M_{2^{s}+1}=\cdots=M_{2^{s+1}}=G^{2^{s+1}},$ and hence $P(s+1)$ holds.
\end{proof}

\end{document}
